%%%%%%%%%%%%%%%%%%%%%%%%%%%%%%%%%%%%%%%%%%%%%%%%%%%%%%%%%%%%%%%%%%%%%%%%%%%%%%%%%%%%%%%%%%%%%%%%
%% Introduction --- reframed draft, 28 November 2026.
%% The framing is the one Munif set: the VISAPP 2027 result is published and settled; the gap is
%% whether corpora of the SAME FORMAT carry it; the hypothesis is that they do, and that on some
%% tasks the encoded image is EQUIVALENT to the bare trace; the method is the statistics that can
%% tell those two apart. The acquisition floor is kept, but as an instrument, not as the headline.
%%
%% Numbers appearing here and their source:
%%   0.645 / 0.614 ............ VISAPP 2027 camera-ready, T2-T4, LPQ+SVM, participant endpoint
%%   0.97 / 0.93 / 0.64 ....... runs/*/d1_capture_floor (ParkinsonHW hw_dataset, Drotar, PaHaW)
%%   1/15 ..................... protocols/PROTOCOL-common.md, addendum H, declared before M1 ran
%%   0.03 ..................... runs/*/g1_split_noise, addendum G; the project's measured band
%%   +0.014 / +0.086 / +0.109 . runs/across_cohorts/per_cohort.csv, addendum F
%%   27 / 19 / 8 / 7 / 6 / 30 . reference/LEVANTAMENTO_DATASETS.md and sec-corpora.tex
%%%%%%%%%%%%%%%%%%%%%%%%%%%%%%%%%%%%%%%%%%%%%%%%%%%%%%%%%%%%%%%%%%%%%%%%%%%%%%%%%%%%%%%%%%%%%%%%

\section{Introduction}
\label{sec:intro}

A trajectory sampled by a digitising tablet can be turned into a picture and handed to a vision
model. In handwriting the construction is unusually direct, because the series already carries a
faithful two-dimensional embedding: the pen path is drawn as ink in the plane in which it was
written, and whatever else the digitiser recorded at each sample --- speed, pressure, azimuth,
altitude --- is painted onto the colour of the stroke at the point where it occurred. Nothing is
averaged away. A vector of kinematic summary statistics, the dominant representation in this
literature \citep{drotar2016pahaw,drotar2020dysgraphia}, records how fast the pen moved; an image
of this kind records \emph{where along the trace} it moved that fast, and leaves the reading of
that spatial arrangement to a texture descriptor or a convolutional network.

We reported one such encoding --- a \emph{pseudodynamic minimap} --- on the PaHaW corpus of
Parkinsonian handwriting, where choosing the painted channel inside each outer training fold
reached $0.645$ participant-level macro F1 against $0.614$ for the bare trace under the same
descriptor and the same classifier \citep{prior2027minimaps}. That result is peer reviewed and
published, and it is also precisely as narrow as its design allows: one corpus, one tablet, one
disease, seventy-five participants, and not a single contrast surviving correction for
multiplicity. Its reviewers asked for independent cohorts; its own future-work section promised
them. This paper is that test, and it is a test of the representation rather than a new one.

\paragraph{The gap.} Results for encodings of this family are, with few exceptions, single-corpus
results: one acquisition device, one recruitment route, and a baseline chosen by the same authors
who chose the encoding. What makes an external test possible here, and unusually cheap, is that
the input to the method is a \emph{format} and not a dataset. Any corpus that stores per-sample pen
coordinates together with a clock can be rendered by identical code, with no redesign and no
re-tuning; the channels a given corpus happens to lack simply remove variants from the menu
(\cref{sec:corpora}). Whether the published result holds elsewhere is therefore an answerable
question, and it had not been answered. It is also an easy question to answer badly. A second
dataset run through a pipeline re-tuned for it measures the analyst, not the representation. What
makes the second corpus evidence about the first is holding the renderer parameters, the
descriptor, the model grid, the split design and the statistics fixed, and writing every
per-corpus decision down before any metric on that corpus is computed.

\paragraph{The question, and the hypothesis.} The useful question is not whether the encoding works
but \emph{where}, and the answer has two halves that need different statistics. The first half is
ordinary: does the margin over the bare trace reproduce outside PaHaW, and how large is it. The
second is the half this literature habitually gets wrong. On which tasks is the encoded image
\emph{as good as} the plain trace? An interval that contains zero is not evidence of equivalence,
and neither is a $p$ above $0.05$: both are also what too few participants look like. Equivalence
is a claim with a threshold, and the threshold has to be fixed before the data are seen for the
claim to mean anything \citep{schuirmann1987,lakens2018equivalence}.

We fix it. Our hypothesis, with its direction declared, is that the encoding pays where the
\emph{task fixes the figure} --- a traced spiral, a copied wave, where patient and control produce
the same shape and only the execution differs --- and is equivalent to the bare trace where the
\emph{produced shape itself carries the label}, because there the geometry already contains what
the painted channel would add. That direction is not a blind prediction, and we do not present it
as one: it is the ordering our own corpus-level analysis reported, restated as a model so that it
can fail. What is untested, and what this paper tests, is whether it survives when the unit of
analysis stops being the corpus.

\paragraph{Why the task is the unit.} Three reasons, and they are why this is a statistics paper
rather than a benchmark table. First, a corpus changes the instrument, the population and the task
all at once, so a contrast between corpora cannot attribute anything to the task property the
hypothesis is about; within a corpus the device and the participants are held fixed by
construction, and the task contrast is the only one that isolates it. Second, a participant
contributes exactly one recording to each task, so a per-task endpoint needs no fusion rule --- the
majority vote over a participant's recordings, which the original protocol required and which mixes
tasks of unequal difficulty, simply disappears. Third, and decisively for this journal: six corpora
are six units. The exact permutation test of a binary task property over six units, two of them on
one side, has a smallest attainable $p$ of $1/15 \approx 0.067$ --- a ceiling we declared before
running that analysis, and one that no amount of signal can pass. The thirty task-level units of
\cref{sec:corpora}, fifteen on each side of the property, two corpora of which contain both kinds
internally, are not bounded that way. They are also not independent, since thirty tasks are
measured on six groups of people, which is what the hierarchical model of \cref{sec:protocol} is
for.

\paragraph{What could make all of this meaningless.} Clinical handwriting corpora are assembled
over months, with patients and controls often recruited through different routes and recorded in
different sessions. Any quantity that drifts with the session --- the sampling interval, the moment
the clock was started --- is then correlated with the label without passing through the disease.
Given those two numbers and no ink whatsoever, a classifier reaches $0.97$ participant-level macro
F1 inside one of these public clinical benchmarks, $0.93$ inside a second, and $0.64$ on PaHaW,
where it ties every pipeline we or others have reported on that corpus. We therefore report each
corpus with that floor beside it, and read each margin against it rather than against zero. The
claim is not that published models read the clock --- our renderers never see an absolute timestamp
--- but that a macro F1 of $0.93$ on a benchmark whose floor is $0.93$ is consistent with having
measured nothing about handwriting at all, and that the reader currently has no way to tell.

\paragraph{Contributions.}
\begin{itemize}
\itemsep0em

\item \textbf{A transfer study defined by the data format, with an auditable screening funnel.}
Twenty-seven candidate corpora were screened against one criterion --- per-sample pen coordinates
with a clock, publicly obtainable without writing to an author. Nineteen fail it, and we give the
reason for each; eight were obtained, seven meet the format contract, six carry a participant
endpoint, and together they contribute thirty task-level units (\cref{sec:corpora}).

\item \textbf{One pipeline, frozen across all of them.} Identical rendering parameters, texture
descriptor, backbone, model grid, nested cross-validation, fusion rule and statistics in every
corpus, with every per-corpus decision --- inclusion rule, task definition, contact rule, unit
conversion, split design --- registered in a protocol document before any classifier was fitted on
that corpus (\cref{sec:protocol}).

\item \textbf{Equivalence tested, not inferred from a null result.} Each task receives one of three
outcomes declared in advance --- superior, equivalent, or inconclusive --- against a margin of
$0.03$ macro F1 that is measured rather than chosen, being this study's own split-to-split
variation. Inconclusive is reported as inconclusive, and the minimum detectable effect is stated
for every unit (\cref{sec:protocol}).

\item \textbf{A descriptor-matched baseline, which removes most of the reported margin.} Giving
each descriptor its own bare-trace baseline shows that the headline margins of this recipe are in
good part a change of descriptor and not of encoding: under a frozen ImageNet backbone no corpus
here gains more than $+0.014$, while under the texture descriptor the two transferring corpora gain
$+0.086$ and $+0.109$ (\cref{sec:results}).

\item \textbf{Mechanism controls that locate the effect, including one that is new.} Shifting the
painted signal along each stroke and permuting it within each stroke destroy only \emph{where} the
signal sits; painting a channel that carries nothing about the writer --- acquisition order, a
smoothness-matched AR(1), white noise --- destroys only \emph{what} the channel says. The first
returns the gain to its baseline; the second does not reproduce it (\cref{sec:results}).

\item \textbf{The acquisition floor, as an instrument rather than a headline.} We define it, give
the procedure, and report every result in the paper as a margin over it (\cref{sec:floor}).

\item \textbf{Failures reported as failures.} A within-corpus test of the hypothesis, registered
with its direction before its splits were drawn, was not confirmed; it is reported as such, and
what it can and cannot mean is bounded by the floor (\cref{sec:stats-moderator}).
\end{itemize}

\subsection{Related work}
\label{sec:related}

\paragraph{Time series as images.} Encoding a sequence as a picture for an image model spans
recurrence plots \citep{eckmann1987recurrence}, Gramian angular fields and Markov transition fields
\citep{wang2015gaf}, and more recently direct rasterisations of trajectories. The handwriting case
differs in that the series already has a faithful two-dimensional embedding, so the encoding does
not have to invent a geometry: the only question is what to do with the remaining channels.

\paragraph{Handwriting as a clinical signal.} Digitised handwriting has been used as a marker for
Parkinson's disease, for developmental dysgraphia and for cognitive decline
\citep{werner2006handwriting}. The dominant representation is a vector of kinematic and pressure
statistics \citep{drotar2016pahaw,drotar2020dysgraphia}; image-based alternatives are more recent,
and our baselines include both.

\paragraph{Replication, equivalence and multiplicity.} Treating a published effect as a hypothesis
to be tested on fresh data, rather than as a result to be extended, is standard practice in fields
that have examined their own reproducibility \citep{osc2015}, and it dictates the shape of this
study: the prior result is the hypothesis, the five new corpora are the test, and the analysis plan
was fixed first. Declaring equivalence requires its own apparatus --- two one-sided tests against a
margin fixed in advance \citep{schuirmann1987,lakens2018equivalence} --- and pooling heterogeneous
units requires a model that admits they are heterogeneous
\citep{higgins2002,hartung2001,viechtbauer2005}. We use all three, and state the limits each of
them has at this sample size.

\paragraph{Shortcut learning and benchmark validity.} That models latch onto acquisition artefacts
rather than the phenomenon is well documented \citep{geirhos2020shortcut}, and in medical imaging
in particular, where scanner and site identity leak into diagnostic labels
\citep{degrave2021shortcut}. Our contribution is not the observation but the instrument: a floor
cheap enough to report with every result, and which we show is not a corner case in this domain.
